\scanpage{344}
\mspage{83}
% archive sheet number: 343

\uline{8}. \uline{Le lemme} d'\uline{homotopie pour les morphismes d'intervalles}

C'est le suivant

\uline{Lemme 1} Soient $\tau,\tau'\in\mathrm{Ty}$, et $f,g:\tau\rightrightarrows\tau'$ deux morphismes de $\underline{\mathrm{Ty}}$. Alors $f,g$ sont homotopes dans la catégorie $\underline{\mathrm{Hom}}\,\underline{\mathrm{int}}(I_\tau,I_{\tau'})$ des morphismes (croissants) d'intervalles orientés de $I_\tau$ dans $I_{\tau'}$.

Cette catégorie pourrait \unsure{licitement|librement} se désigner par $\underline{\mathrm{Hom}}_{\underline{\mathrm{Ty}}}(\tau,\tau')$, ou $\underline{\mathrm{Hom}}(\tau,\tau')$. C'est une sous-catégorie \deleted{pleine} de la catégorie $\underline{\mathrm{Hom}}(I_\tau,I_{\tau'})$ de \uline{tous} les foncteurs de $I_\tau$ dans $I_{\tau'}$, et une sous-catégorie pleine de la catégorie $\underline{\mathrm{Hom}}^{\unsure{\varepsilon|E}}(I_\tau,I_{\tau'})$ des foncteurs \inserted{$f$} qui sont compatibles avec les données de \unsure{sources|source} et de \unsure{buts|but}, i.e. qui rendent \unsure{commutatif|commutatifs} le diagramme
\diagram{\begin{tikzcd}[ampersand replacement=\&,column sep=small]
\&\varepsilon=\{0,1\}\arrow[dl,hook,"{(\sigma_\tau,\beta_\tau)}"']\arrow[dr,hook,"{(\sigma_{\tau'},\beta_{\tau'})}"]\&\\
I_\tau\arrow[rr,"f"]\&\&I_{\tau'}
\end{tikzcd}}

Si $f,g$ \deleted{\unsure{ont}} sont dans cette catégorie, \deleted{alors} (qui soient morphismes croissants d'ailleurs), une flèche de $f$ dans $g$ est donc \deleted{\unsure{un hom}} la donnée, pour tout $i\in I_\tau$, d'une \deleted{\unsure{hom}} \inserted{flèche} $f(i)\xrightarrow{\unsure{w|u}_i}g(i)$ de $I_{\tau'}$, \unsure{que|pour|sous} les deux conditions a) si $i\in\{0_\tau,1_\tau\}$, $\unsure{w'|w^{i}}=\mathrm{id}$ et b) le carré commutatif habituel pour une flèche $i\to j$ de $I_\tau$. (\unsure{Il} suffit d'avoir cette compatibilité pour les flèches de transition $\underline{u}_i$\inserted{$:i-1\to i$} de $I_\tau$ ($i\in[1,\mathrm{long}\,\tau]$

Une autre façon de décrire cette catégorie \inserted{$\underline{\mathrm{Hom}}^{\varepsilon}$}, c'est de \margin{\unsure{$\frac{=}{i}$}} dire que c'est \deleted{celle des} la catégorie \deleted{des $\underline{\mathrm{Ch}}_\tau(I_{\tau'};0_{\tau'},1_{\tau'})$}
